\ssscp{Shifting a core cultural metaphor: ‘seat of emotions, character’}{02-03-01-06}
Just as structural change indicates intense contact over long
periods of time, so also one would not expect entrenched \it{core cultural metaphors}
relating to daily relationships (such as \ili{English} \it{his heart
was broken}) to shift easily (see \citealp{LakoffJohnson1980}).
In this case it is also significant that so many languages shift
to the same metaphor,
but sporadically, and in languages spread over hundreds
and even thousands of kilometers.
Yet that is what we are seeing in the region,
as has been noted in \citet[148--150]{DonohueGrimes2008}, and \citet{GrimesCIP}.

Many languages around the world use a \it{figurative sense} associated
with a body organ (e.g.
kidney, heart, liver, stomach) to indicate something like ‘seat
of emotions, character and social relations’.
This is seen in Indo-European languages,
and it is this figurative sense that is involved in much poetry,
literature, and songs -- not the literal sense of the body part.
Many Indo-European languages use a figurative sense of \it{heart}
‘1) organ that pumps blood,
2) seat of emotions
and character’, as illustrated in \qf{02-36}.

{\wg\tblr{>{\hspace{-\tabcolsep}}l<{\hspace{-\tabcolsep}}>{\hspace{-\tabcolsep}\enspace}llll}
\lsptoprule
\tbx{02-36}		&	Proto-Indo-European	&	*ḱérd				&	/kʸerd/	&	Gloss		\\
\midrule
&	\ili{English}			&	{\hr}heart					&	/hɑːt/			&	1) heart; 2) emotions		\\
&	Icelandic		&	{\hr}hjarta					&	/hyar̥ta/		&	1) heart; 2) emotions		\\
&	Irish				&	{\hr}croí						&	/kɾˠiː/			&	1) heart; 2) emotions		\\
&	Greek				&	{\hr}καρδία 				&	/karðia/		&	1) heart; 2) emotions		\\
&	Armenian		&	{\hr}\Armenian{սիրտ}&	/sirt/			&	1) heart; 2) emotions		\\
&	Russian			&	{\hr}сердце 				&	/sʸerʦe/		&	1) heart; 2) emotions		\\
&	Latvian			&	{\hr}sirds	 				&	/sirts/			&	1) heart; 2) emotions		\\
&	\ili{Sanskrit}		&	{\hr}\\ili{Sanskrit}{हृदय}	&	/ɦr̩d̪aya/		&	1) heart; 2) emotions		\\
&	Spanish			&	{\hr}corazón				&	/koɾaˈson/,	&	1) heart; 2) emotions	\\ \hhline{~}
&							&											&	/koɾaˈθon/	&			\\
&	Portuguese	&	{\hr}coração				&	/kuɾasãũ/		&	1) heart; 2) emotions		\\ \hhline{~}
&							&											&	/koɾasãũ/		&				\\
&	French			&	{\hr}coeur					&	/kœʁ/				&	1) heart; 2) emotions		\\
\lspbottomrule
\end{tabular}\mbox{}\\}

Many western Austronesian languages use a figurative sense of
PMP *qatay ‘1) liver’ (the organ that cleans blood) to also indicate
‘2) seat of emotions,
character and social relations’.
This is also the stuff of poetry,
literature, and songs
(data in \qf{02-37} and \qf{02-38} are gleaned from
\citealp{Tryon1995} and \citealp{BlustTrussel2020}).

{\wg\tblr{>{\hspace{-\tabcolsep}}l<{\hspace{-\tabcolsep}}>{\hspace{-\tabcolsep}\enspace}lll}
\lsptoprule
\tbx{02-37}&	\ili{Tagalog}	&	atay	&	1) liver; 2) emotions		\\
&	\ili{Malay}	&	hati	&	1) liver; 2) emotions		\\
&	Sunda	&	hate	&	1) liver; 2) emotions		\\
&	\ili{Makasar}	&	ate	&	1) liver; 2) emotions		\\
&	\ili{Wolio}	&	ate	&	1) liver; 2) emotions		\\
&	\ili{Sasak}	&	ate	&	1) liver; 2) emotions		\\
\lspbottomrule
\end{tabular}\mbox{}\\}

Some Austronesian languages \ili{east} of the Blust line (mostly \tsc{\ili{Bima}-Lembata}
languages) also maintain this figurative sense of ‘liver’.

{\wg\tblr{>{\hspace{-\tabcolsep}}l<{\hspace{-\tabcolsep}}>{\hspace{-\tabcolsep}\enspace}lll}
\lsptoprule
\tbx{02-38}	&	\ili{Bima}	&	ade	&	1) liver; 2) emotions		\\
	&	\ili{Ngadha}	&	ate	&	1) liver; 2) emotions		\\
	&	\ili{Sika}	&	ʔβate	&	1) liver; 2) emotions		\\
	&	\ili{Kambera}	&	eti	&	1) liver; 2) emotions		\\
\lspbottomrule
\end{tabular}\mbox{}\\}

Yet throughout our target region we find that many Austronesian
languages (some of which also maintain a reflex of PMP *qatay
‘liver’ for the organ ‘liver’ or ‘spleen’),
have shifted to use a reflex of PMP *daləm ‘inside’ for the figurative
sense of ‘seat of emotions,
character and social relations’.
This is the same word used in locative phrases such as ‘inside
the house’, but also to refer to ‘his/her insides/character/seat of emotions’.
Example \qf{02-39} shows a small sample of some of languages that use
‘inside’ figuratively in this way.

{\wg\tblr{>{\hspace{-\tabcolsep}}l<{\hspace{-\tabcolsep}}>{\hspace{-\tabcolsep}\enspace}llll}
\lsptoprule
\tbx{02-39}	&	\ili{Dhao}\footnotemark[20]
	&	ɗara	&	1) inside, 2) emotions	&	\tsc{\ili{Bima}-Lembata}		\\
%	&	\ili{Havu}\footnote{}	&	ɗara	&	1) inside, 2) emotions	&	\tsc{\ili{Bima}-Lembata}		\\
	&	\ili{Dela}	&	rala-n	&	1) inside, 2) emotions	&	\tsc{Timor-Babar}		\\
	&	\ili{Tii}	&	dale-n	&	1) inside, 2) emotions	&	\tsc{Timor-Babar}		\\
	&	\ili{Helong}	&	dale-n	&	1) inside, 2) emotions	&	\tsc{Timor-Babar}		\\
	&	\ili{Tetun}\footnotemark[21]
	&	lara-n	&	1) inside, 2) emotions
	&	\tsc{Timor-Babar}		\\
	&	\ili{Luang}	&	ralma ni	&	1) inside, 2) emotions	&	\tsc{Timor-Babar}		\\
	&	S. \ili{Mambae}	&	hua-lala	&	‘fruit insides/heart insides’	&	\tsc{\ili{Central} Timor}		\\
	&	\ili{Kemak}	&	lara-n	&	1) inside, 2) emotions	&	\tsc{\ili{Central} Timor}		\\
	&	\ili{Buru}	&	lale-n	&	1) inside, 2) emotions	&	\tsc{Sula-Buru}		\\
	&	\ili{Asilulu}	&	lale-	&	1) inside, 2) emotions	&	\tsc{\ili{Ambon}-Seram}		\\
	&	\ili{Yamdena}	&	dalamy	&	1) inside, 2) emotions	&	\tsc{STB}		\\
\lspbottomrule
\end{tabular}\mbox{}\\}

\footnotetext[20]{
		While \ili{Dhao} consistently uses \it{ɗara} for ‘seat of emotions’,
		closely related \ili{Havu} is split, using both \it{ade} ‘liver’
		and \it{ɗara} ‘insides’ for ‘seat of emotions’,
		with \it{ade} being the more frequent.
		So \ili{Havu} has \it{adu ade} ‘resistant (lit: hard liver)’,
		and \it{meŋəlu ɗara} ‘happy (lit: cool insides)’,
		and many other similar expressions.
		Some expressions can use either,
		such as \it{pe-rui ade} or \it{pe-rui ɗara}
		‘encourage (lit: \tsc{cause}-strong liver/insides)’.
		So while \ili{Dhao} has shifted to ‘insides’,
		its nearest linguistic relative straddles the fence,
		but remains mostly in the ‘liver’ camp.}\addtocounter{footnote}{1}
\footnotetext[21]{Grimes has worked extensively with both \ili{Tetun} Dili [tdt]
		and \ili{Tetun} [tet], and it is clear that the gloss in
		\citet{BlustTrussel2020} for \ili{Tetun} Dili \it{laran} glossed merely as
		‘in, inside’ while not wrong,
		is also incomplete, as reading newspapers or listening to love
		songs played around Dili reveals daily.
		See also B. \citet{Grimes2018}.
		The gloss for \ili{Havu} \it{ɗara} in \citet{BlustTrussel2020} as ‘in,
		within (both spatial and temporal)’ is similarly incomplete (see
		\cite[254]{Grimes2008,Grimes2010a}).
		This discussion of a secondary sense of words in comparative
		studies perhaps illustrates some of the limitations of working
		with word lists, simple glossaries, or simple glosses in a linguistic paper,
		and emphasises the need for serious corpus-based dictionaries
		that attend to polysemy.
	}\addtocounter{footnote}{1}

%\begin{envfn}
%\footnotetext{
		%While \ili{Dhao} consistently uses \it{ɗara} for ‘seat of emotions’,
		%closely related \ili{Havu} is split, using both \it{ade} ‘liver’
		%and \it{ɗara} ‘insides’ for ‘seat of emotions’,
		%with \it{ade} being the more frequent.
		%So \ili{Havu} has \it{adu ade} ‘resistant (lit: hard liver)’,
		%and \it{meŋəlu ɗara} ‘happy (lit: cool insides)’,
		%and many other similar expressions.
		%Some expressions can use either,
		%such as \it{pe-rui ade} or \it{pe-rui ɗara}
		%‘encourage (lit: \tsc{cause}-strong liver/insides)’.
		%So while \ili{Dhao} has shifted to ‘insides’,
		%its nearest linguistic relative straddles the fence,
		%but remains mostly in the ‘liver’ camp.}
%\end{envfn}
%\begin{envfn}
%\footnotetext{Grimes has worked extensively with both \ili{Tetun} Dili [tdt]
		%and \ili{Tetun} [tet], and it is clear that the gloss in
		%\citet{BlustTrussel2020} for \ili{Tetun} Dili \it{laran} glossed merely as
		%‘in, inside’ while not wrong,
		%is also incomplete, as reading newspapers or listening to love
		%songs played around Dili reveals daily.
		%See also B. \citet{Grimes2018}.
		%The gloss for \ili{Havu} \it{ɗara} in \citet{BlustTrussel2020} as ‘in,
		%within (both spatial and temporal)’ is similarly incomplete (see
		%\cite[254]{Grimes2008}, \cite{Grimes2010a}).
		%This discussion of a secondary sense of words in comparative
		%studies perhaps illustrates some of the limitations of working
		%with word lists, simple glossaries, or simple glosses in a linguistic paper,
		%and emphasises the need for serious corpus-based dictionaries
		%that attend to polysemy.
	%}
%\end{envfn}

This naturally raises the question as to what has brought about
such a monumental reorientation from figurative senses of *qatay
‘liver’ to figurative senses of *daləm ‘inside’ for emotions
and social relations? Since we find this same metaphor in Papuan
languages, it is not unreasonable to once again infer that this shift reflects
contact-induced change in these Austronesian languages to model
on the metaphors of Papuan languages,
as shown in \qf{02-40}.

{\wg\tblr{>{\hspace{-\tabcolsep}}l<{\hspace{-\tabcolsep}}>{\hspace{-\tabcolsep}\enspace}llll}
\lsptoprule
\tbx{02-40}	&	\ili{Teiwa}	&	gom	&	1) inside; 2) heart	&	TAP		\\
	&	Pura	&	omi	&	1) inside; 2) place of belief	&	TAP		\\
	&	Abolo	&	ʊm	&	1) inside; 2) emotions	&	TAP		\\
	&	\ili{Wersing}	&	mira	&	1) inside; 2) emotions	&	TAP		\\
\lspbottomrule
\end{tabular}\mbox{}\\}

For such a fundamental worldview concept to shift to a different
metaphor,\footnote{
		“The power to reason about so abstract an
		idea as life comes very largely through metaphor.”
		\citep[62]{LakoffTurner1989}.}
once again it implies there has been intense
contact, bilingualism, and accommodation over a great period of time.
And again, this does not just happen in a vacuum.
The traces of this contact-induced change are seen throughout
the region, both in proximity to
and far from where any Papuan languages are spoken today.
A perusal of PMP *daləm ‘inside’ in \citet{BlustTrussel2020},
shows that the figurative sense of ‘seat
of emotions’ is recorded only occasionally west of the Blust
line, but is pervasive in the contact zone of Wallacea
and extending throughout Oceanic languages.